% =========================================================================
% Materials and methods -- fragment for the PLoS Computational Biology
% template (% =========================================================================
% Materials and methods -- fragment for the PLoS Computational Biology
% template (% =========================================================================
% Materials and methods -- fragment for the PLoS Computational Biology
% template (% =========================================================================
% Materials and methods -- fragment for the PLoS Computational Biology
% template (\input{methods.tex}). Requires amsmath (loaded by the template).
% Citation keys are placeholders to be resolved in the project .bib file.
% =========================================================================

\section*{Materials and methods}

Our experimental logic proceeds in four steps. First, we trained recurrent
neural networks (RNNs) exclusively to \emph{decode} (denoise) noisy
observations of a temporally structured stimulus stream---Bach
compositions---without any explicit predictive objective. Second, we asked
whether \emph{prediction emerges} as a by-product of this training: freezing
the network, we probed its hidden state with a linear readout trained to
report the upcoming stimulus. Third, we \emph{benchmarked} the probed
prediction accuracy against a ladder of reference models whose expected
errors are ordered by construction, allowing us to attribute performance to
specific computational capacities. Fourth, we tested whether \emph{prediction
errors are encoded} in the population activity of the trained networks, as
posited by predictive-coding accounts of cortical
processing~\cite{rao1999predictive,friston2005theory}.

\subsection*{Stimulus corpus}

We assembled a corpus of works by J.~S.~Bach from publicly available MIDI
archives.\footnote{jsbach.net, metronimo.com, piano-midi.de, sound.jp/bach-ken,
dardel.info, classicalarchives.com, and suzumidi.com.} Each MIDI file was
parsed into an event-based sequence of chords: a new element was created
whenever the set of simultaneously sounding pitches changed, so that the
sequence index $t$ advances with the musical surface rather than with clock
time. This discretisation removes tempo and rhythm while preserving the
harmonic and voice-leading structure of the music. Pitches were collapsed
into the 12 pitch classes, yielding for each time step a multi-hot binary
vector $\mathbf{y}_t \in \{0,1\}^{12}$ whose $k$-th component indicates
whether pitch class $k$ is sounding (mean polyphony $\approx 2.3$
simultaneous pitch classes). Files belonging to the same composition
(different arrangements or movements sharing a BWV catalogue number) were
concatenated, separated by brief silences.

The corpus comprised 700 compositions, split at the composition level into
training (70\%, 490 works), validation (10\%, 70 works) and test sets (20\%,
140 works). Splitting by BWV number guarantees that no arrangement or
movement of a test composition was ever seen during training.

\subsection*{Noise model and decoding task}

Networks never observed the stimulus directly. At each time step the input
was a noise-corrupted observation
\begin{equation}
  \mathbf{x}_t \;=\; \mathbf{y}_t + \sigma\,\boldsymbol{\varepsilon}_t,
  \qquad \boldsymbol{\varepsilon}_t \sim \mathcal{N}(\mathbf{0}, I_{12}),
  \label{eq:noise}
\end{equation}
with the noise amplitude $\sigma$ treated as an experimental factor spanning
nine levels: a near-noiseless control ($\sigma = 0.001$) and $\sigma \in
\{0.2, 0.4, 0.6, 0.8, 1.0, 1.2, 1.6, 2.0\}$, sampled densely in the range
where the single-observation discriminability $d' = 1/\sigma$ crosses unity
and sparsely near the two asymptotic regimes. At the upper end of this range
the single-step signal-to-noise ratio is far below unity, so that accurate
decoding is impossible without integrating information over time---the
regime in which an internal model of the stimulus dynamics becomes
computationally necessary, in analogy to optimal filtering of a hidden
dynamical process~\cite{kalman1960}. During training,
an independent noise realisation was drawn at every presentation of a
stimulus segment, making the effective training set unbounded and preventing
memorisation of any particular noise pattern. Validation and test sets used
fixed noise realisations, making model selection and evaluation
deterministic.

\subsection*{Network architecture}

The core model was a single-layer gated recurrent unit network
(GRU~\cite{cho2014gru}) with $n$ hidden units, $n \in \{8, 16, 32, 64, 128,
256\}$; the smallest size exceeds the stimulus dimensionality, so no network
was representationally bottlenecked below its input. Two parallel readout heads, each a linear map followed by an
element-wise sigmoid, projected the hidden state $\mathbf{h}_t \in
\mathbb{R}^{n}$ to the stimulus space: an \emph{observation head}
$\hat{\mathbf{y}}_t = \varsigma(W_{o}\mathbf{h}_t + \mathbf{b}_o)$ reporting
the current chord, and a \emph{prediction head} $\mathbf{p}_t =
\varsigma(W_{p}\mathbf{h}_t + \mathbf{b}_p)$ reporting the upcoming chord.
Outputs were treated as independent Bernoulli probabilities over the 12
pitch classes and scored with the binary cross-entropy (BCE),
\begin{equation}
  \mathcal{L}_{\mathrm{obs}}
    = \big\langle \mathrm{BCE}(\hat{\mathbf{y}}_t, \mathbf{y}_t) \big\rangle_t,
  \qquad
  \mathcal{L}_{\mathrm{pred}}
    = \big\langle \mathrm{BCE}(\mathbf{p}_t, \mathbf{y}_{t+1}) \big\rangle_t,
  \label{eq:objectives}
\end{equation}
where $\langle\cdot\rangle_t$ averages over time steps and pitch classes.
The hidden state was initialised to zero at the beginning of each sequence.

\subsection*{Stage 1: training to decode}

In the first training stage only the recurrent weights and the observation
head were optimised, under the decoding objective
$\mathcal{L}_{\mathrm{obs}}$ alone; the prediction head was frozen at its
random initialisation and contributed no gradient. Networks were trained
with Adam~\cite{kingma2015adam} (learning rate $0.02$) on mini-batches of
512 segments of 512 consecutive time steps, drawn by advancing through
randomly selected training compositions. Each network received a fixed
budget of 20{,}000 batches; no early stopping was used, so that every cell
of the $\sigma \times n$ design received an identical amount of training.
Decoding error on the validation set was measured every 100 batches, and the
network state with the lowest validation error was retained as the final
model (best-checkpoint selection). Five networks with independent random
initialisations were trained for each of the $9 \times 6$ combinations of
noise amplitude and network size.

\subsection*{Stage 2: probing for emergent prediction}

To ask whether the decoding-trained networks had learned to anticipate their
input, we froze the recurrent weights and the observation head and trained
\emph{only} the linear prediction head, under
$\mathcal{L}_{\mathrm{pred}}$, for 3{,}000 batches with otherwise identical
hyperparameters. Because the recurrent dynamics were fixed, this readout
cannot create new temporal computations: its accuracy is an estimate of the
predictive information already present in the hidden state. Any prediction
performance exceeding that of the memoryless reference models below must
therefore have emerged spontaneously from the denoising objective of
Stage~1.

\subsection*{Reference models}

Prediction accuracy is only interpretable relative to the intrinsic
predictability of the stimulus. We therefore constructed a ladder of
reference models whose expected losses, at their respective optima,
correspond to conditional entropies of the stimulus process and are
provably ordered. Writing $H(\cdot\,|\,\cdot)$ for the conditional entropy
of a pitch class, independence of the observation noise implies the Markov
chain $\mathbf{x}_t \to \mathbf{y}_t \to \mathbf{y}_{t+1}$, and hence
\begin{equation}
  \underbrace{H(\mathbf{y}_{t+1} \,|\, \mathbf{y}_t)}_{\text{Markov-low}}
  \;\le\;
  \underbrace{H(\mathbf{y}_{t+1} \,|\, \mathbf{x}_t)}_{\text{Markov-high}}
  \;\le\;
  \underbrace{H(\mathbf{y}_{t+1})}_{\text{global distribution}} .
  \label{eq:ladder}
\end{equation}
The corresponding estimators were:
\begin{description}
  \item[Global distribution.] The optimal static predictor under the BCE
  loss: the vector of per-pitch-class marginal frequencies estimated from
  the training set, constantly predicted at every time step. This model has
  no access to temporal structure and upper-bounds the error of any
  informed predictor.
  \item[Markov-high.] A feedforward network (one hidden layer of 256 units)
  mapping the noisy observation $\mathbf{x}_t$ to $\mathbf{y}_{t+1}$: the
  best one-step predictor operating, like the RNN, through the noisy
  channel, but without memory.
  \item[Markov-low.] The same architecture mapping the \emph{clean} chord
  $\mathbf{y}_t$ to $\mathbf{y}_{t+1}$, trained and evaluated with clean
  input: the best one-step predictor with perfect knowledge of the current
  state. Its error is independent of $\sigma$ and lower-bounds all
  first-order predictors.
  \item[Prediction-trained GRU.] A single-layer GRU of 256 units trained
  directly on $\mathcal{L}_{\mathrm{pred}}$, providing an empirical
  estimate of the best prediction attainable by the architecture itself.
  \item[Memoryless decoder.] For the decoding task, a feedforward network
  (three hidden layers of 256 units) mapping $\mathbf{x}_t$ to
  $\mathbf{y}_t$, quantifying how much decoding accuracy is attainable
  without temporal integration.
\end{description}
Because the bound interpretation of Eq~\eqref{eq:ladder} requires reference
models operating near their optima, all reference models were
capacity-saturated (256 hidden units, the maximum of the network-size
sweep) and computed once per noise level, rather than width-matched to each
RNN of the sweep. They were trained with the same optimiser, batch
geometry, validation cadence and checkpoint selection as the core networks,
with a budget of 25{,}000 batches.

\subsection*{Evaluation and statistics}

All models were evaluated on the 140 held-out test compositions, presented
in full and never seen during training or model selection. The test-set
noise realisation was generated once from a fixed seed and shared by all
models and all noise amplitudes (common random numbers: the same
$\boldsymbol{\varepsilon}$ scaled by $\sigma$), so that comparisons between
models are paired both by composition and by noise realisation. Performance
was summarised as the mean BCE per composition; for the core networks,
per-composition errors were first averaged across the five training runs, so
that the composition is the unit of analysis throughout. This yields for
each model a distribution of 140 paired errors. Pairwise comparisons
between models used
two-sided Wilcoxon signed-rank tests across test compositions, with effect
sizes reported as Cohen's $d$ computed with the pooled standard deviation.
For summary maps spanning the full design, $p$-values were
Bonferroni-corrected for the $9 \times 6 = 54$ cells of the sweep.

\subsection*{Encoding of prediction errors}

Finally, we asked whether the trained networks explicitly represent their
own prediction errors, defined at each time step as the deviation of the
incoming observation from the network's prediction,
\begin{equation}
  \mathbf{e}_t \;=\; \mathbf{x}_{t+1} - \mathbf{p}_t .
  \label{eq:pe}
\end{equation}
Note that $\mathbf{e}_t$ is defined with respect to the noisy observation,
the only signal available to the network, so its magnitude scales with
$\sigma$. We related the population prediction error $\|\mathbf{e}_t\|^2$ to
two signatures of the hidden dynamics: the state energy
$\|\mathbf{h}_{t+1}\|^2$ and the state-change energy $\|\mathbf{h}_{t+1} -
\mathbf{h}_t\|^2$ (all quantities averaged over their components). For each
test composition we computed the Pearson correlation between the
prediction-error time course and each signature. To test whether the full
error vector, rather than its magnitude alone, is linearly decodable from
the population, we fitted ordinary least-squares regressions from
$\mathbf{h}_{t+1}$ (and from $\mathbf{h}_{t+1}-\mathbf{h}_t$) to
$\mathbf{e}_t$ on the training compositions and evaluated the mean squared
decoding error on each test composition, against a baseline decoder that
constantly predicts the training-set mean error vector.

\subsection*{Implementation}

All models were implemented in PyTorch~\cite{paszke2019pytorch}; MIDI
parsing used mido, and statistical analyses used SciPy and
scikit-learn~\cite{pedregosa2011sklearn}. The full pipeline---corpus
retrieval, model training, benchmarking and analysis---is available at
\url{https://github.com/USERNAME/BayesPlusBach}. % TODO: final repo URL
). Requires amsmath (loaded by the template).
% Citation keys are placeholders to be resolved in the project .bib file.
% =========================================================================

\section*{Materials and methods}

Our experimental logic proceeds in four steps. First, we trained recurrent
neural networks (RNNs) exclusively to \emph{decode} (denoise) noisy
observations of a temporally structured stimulus stream---Bach
compositions---without any explicit predictive objective. Second, we asked
whether \emph{prediction emerges} as a by-product of this training: freezing
the network, we probed its hidden state with a linear readout trained to
report the upcoming stimulus. Third, we \emph{benchmarked} the probed
prediction accuracy against a ladder of reference models whose expected
errors are ordered by construction, allowing us to attribute performance to
specific computational capacities. Fourth, we tested whether \emph{prediction
errors are encoded} in the population activity of the trained networks, as
posited by predictive-coding accounts of cortical
processing~\cite{rao1999predictive,friston2005theory}.

\subsection*{Stimulus corpus}

We assembled a corpus of works by J.~S.~Bach from publicly available MIDI
archives.\footnote{jsbach.net, metronimo.com, piano-midi.de, sound.jp/bach-ken,
dardel.info, classicalarchives.com, and suzumidi.com.} Each MIDI file was
parsed into an event-based sequence of chords: a new element was created
whenever the set of simultaneously sounding pitches changed, so that the
sequence index $t$ advances with the musical surface rather than with clock
time. This discretisation removes tempo and rhythm while preserving the
harmonic and voice-leading structure of the music. Pitches were collapsed
into the 12 pitch classes, yielding for each time step a multi-hot binary
vector $\mathbf{y}_t \in \{0,1\}^{12}$ whose $k$-th component indicates
whether pitch class $k$ is sounding (mean polyphony $\approx 2.3$
simultaneous pitch classes). Files belonging to the same composition
(different arrangements or movements sharing a BWV catalogue number) were
concatenated, separated by brief silences.

The corpus comprised 700 compositions, split at the composition level into
training (70\%, 490 works), validation (10\%, 70 works) and test sets (20\%,
140 works). Splitting by BWV number guarantees that no arrangement or
movement of a test composition was ever seen during training.

\subsection*{Noise model and decoding task}

Networks never observed the stimulus directly. At each time step the input
was a noise-corrupted observation
\begin{equation}
  \mathbf{x}_t \;=\; \mathbf{y}_t + \sigma\,\boldsymbol{\varepsilon}_t,
  \qquad \boldsymbol{\varepsilon}_t \sim \mathcal{N}(\mathbf{0}, I_{12}),
  \label{eq:noise}
\end{equation}
with the noise amplitude $\sigma$ treated as an experimental factor spanning
nine levels: a near-noiseless control ($\sigma = 0.001$) and $\sigma \in
\{0.2, 0.4, 0.6, 0.8, 1.0, 1.2, 1.6, 2.0\}$, sampled densely in the range
where the single-observation discriminability $d' = 1/\sigma$ crosses unity
and sparsely near the two asymptotic regimes. At the upper end of this range
the single-step signal-to-noise ratio is far below unity, so that accurate
decoding is impossible without integrating information over time---the
regime in which an internal model of the stimulus dynamics becomes
computationally necessary, in analogy to optimal filtering of a hidden
dynamical process~\cite{kalman1960}. During training,
an independent noise realisation was drawn at every presentation of a
stimulus segment, making the effective training set unbounded and preventing
memorisation of any particular noise pattern. Validation and test sets used
fixed noise realisations, making model selection and evaluation
deterministic.

\subsection*{Network architecture}

The core model was a single-layer gated recurrent unit network
(GRU~\cite{cho2014gru}) with $n$ hidden units, $n \in \{8, 16, 32, 64, 128,
256\}$; the smallest size exceeds the stimulus dimensionality, so no network
was representationally bottlenecked below its input. Two parallel readout heads, each a linear map followed by an
element-wise sigmoid, projected the hidden state $\mathbf{h}_t \in
\mathbb{R}^{n}$ to the stimulus space: an \emph{observation head}
$\hat{\mathbf{y}}_t = \varsigma(W_{o}\mathbf{h}_t + \mathbf{b}_o)$ reporting
the current chord, and a \emph{prediction head} $\mathbf{p}_t =
\varsigma(W_{p}\mathbf{h}_t + \mathbf{b}_p)$ reporting the upcoming chord.
Outputs were treated as independent Bernoulli probabilities over the 12
pitch classes and scored with the binary cross-entropy (BCE),
\begin{equation}
  \mathcal{L}_{\mathrm{obs}}
    = \big\langle \mathrm{BCE}(\hat{\mathbf{y}}_t, \mathbf{y}_t) \big\rangle_t,
  \qquad
  \mathcal{L}_{\mathrm{pred}}
    = \big\langle \mathrm{BCE}(\mathbf{p}_t, \mathbf{y}_{t+1}) \big\rangle_t,
  \label{eq:objectives}
\end{equation}
where $\langle\cdot\rangle_t$ averages over time steps and pitch classes.
The hidden state was initialised to zero at the beginning of each sequence.

\subsection*{Stage 1: training to decode}

In the first training stage only the recurrent weights and the observation
head were optimised, under the decoding objective
$\mathcal{L}_{\mathrm{obs}}$ alone; the prediction head was frozen at its
random initialisation and contributed no gradient. Networks were trained
with Adam~\cite{kingma2015adam} (learning rate $0.02$) on mini-batches of
512 segments of 512 consecutive time steps, drawn by advancing through
randomly selected training compositions. Each network received a fixed
budget of 20{,}000 batches; no early stopping was used, so that every cell
of the $\sigma \times n$ design received an identical amount of training.
Decoding error on the validation set was measured every 100 batches, and the
network state with the lowest validation error was retained as the final
model (best-checkpoint selection). Five networks with independent random
initialisations were trained for each of the $9 \times 6$ combinations of
noise amplitude and network size.

\subsection*{Stage 2: probing for emergent prediction}

To ask whether the decoding-trained networks had learned to anticipate their
input, we froze the recurrent weights and the observation head and trained
\emph{only} the linear prediction head, under
$\mathcal{L}_{\mathrm{pred}}$, for 3{,}000 batches with otherwise identical
hyperparameters. Because the recurrent dynamics were fixed, this readout
cannot create new temporal computations: its accuracy is an estimate of the
predictive information already present in the hidden state. Any prediction
performance exceeding that of the memoryless reference models below must
therefore have emerged spontaneously from the denoising objective of
Stage~1.

\subsection*{Reference models}

Prediction accuracy is only interpretable relative to the intrinsic
predictability of the stimulus. We therefore constructed a ladder of
reference models whose expected losses, at their respective optima,
correspond to conditional entropies of the stimulus process and are
provably ordered. Writing $H(\cdot\,|\,\cdot)$ for the conditional entropy
of a pitch class, independence of the observation noise implies the Markov
chain $\mathbf{x}_t \to \mathbf{y}_t \to \mathbf{y}_{t+1}$, and hence
\begin{equation}
  \underbrace{H(\mathbf{y}_{t+1} \,|\, \mathbf{y}_t)}_{\text{Markov-low}}
  \;\le\;
  \underbrace{H(\mathbf{y}_{t+1} \,|\, \mathbf{x}_t)}_{\text{Markov-high}}
  \;\le\;
  \underbrace{H(\mathbf{y}_{t+1})}_{\text{global distribution}} .
  \label{eq:ladder}
\end{equation}
The corresponding estimators were:
\begin{description}
  \item[Global distribution.] The optimal static predictor under the BCE
  loss: the vector of per-pitch-class marginal frequencies estimated from
  the training set, constantly predicted at every time step. This model has
  no access to temporal structure and upper-bounds the error of any
  informed predictor.
  \item[Markov-high.] A feedforward network (one hidden layer of 256 units)
  mapping the noisy observation $\mathbf{x}_t$ to $\mathbf{y}_{t+1}$: the
  best one-step predictor operating, like the RNN, through the noisy
  channel, but without memory.
  \item[Markov-low.] The same architecture mapping the \emph{clean} chord
  $\mathbf{y}_t$ to $\mathbf{y}_{t+1}$, trained and evaluated with clean
  input: the best one-step predictor with perfect knowledge of the current
  state. Its error is independent of $\sigma$ and lower-bounds all
  first-order predictors.
  \item[Prediction-trained GRU.] A single-layer GRU of 256 units trained
  directly on $\mathcal{L}_{\mathrm{pred}}$, providing an empirical
  estimate of the best prediction attainable by the architecture itself.
  \item[Memoryless decoder.] For the decoding task, a feedforward network
  (three hidden layers of 256 units) mapping $\mathbf{x}_t$ to
  $\mathbf{y}_t$, quantifying how much decoding accuracy is attainable
  without temporal integration.
\end{description}
Because the bound interpretation of Eq~\eqref{eq:ladder} requires reference
models operating near their optima, all reference models were
capacity-saturated (256 hidden units, the maximum of the network-size
sweep) and computed once per noise level, rather than width-matched to each
RNN of the sweep. They were trained with the same optimiser, batch
geometry, validation cadence and checkpoint selection as the core networks,
with a budget of 25{,}000 batches.

\subsection*{Evaluation and statistics}

All models were evaluated on the 140 held-out test compositions, presented
in full and never seen during training or model selection. The test-set
noise realisation was generated once from a fixed seed and shared by all
models and all noise amplitudes (common random numbers: the same
$\boldsymbol{\varepsilon}$ scaled by $\sigma$), so that comparisons between
models are paired both by composition and by noise realisation. Performance
was summarised as the mean BCE per composition; for the core networks,
per-composition errors were first averaged across the five training runs, so
that the composition is the unit of analysis throughout. This yields for
each model a distribution of 140 paired errors. Pairwise comparisons
between models used
two-sided Wilcoxon signed-rank tests across test compositions, with effect
sizes reported as Cohen's $d$ computed with the pooled standard deviation.
For summary maps spanning the full design, $p$-values were
Bonferroni-corrected for the $9 \times 6 = 54$ cells of the sweep.

\subsection*{Encoding of prediction errors}

Finally, we asked whether the trained networks explicitly represent their
own prediction errors, defined at each time step as the deviation of the
incoming observation from the network's prediction,
\begin{equation}
  \mathbf{e}_t \;=\; \mathbf{x}_{t+1} - \mathbf{p}_t .
  \label{eq:pe}
\end{equation}
Note that $\mathbf{e}_t$ is defined with respect to the noisy observation,
the only signal available to the network, so its magnitude scales with
$\sigma$. We related the population prediction error $\|\mathbf{e}_t\|^2$ to
two signatures of the hidden dynamics: the state energy
$\|\mathbf{h}_{t+1}\|^2$ and the state-change energy $\|\mathbf{h}_{t+1} -
\mathbf{h}_t\|^2$ (all quantities averaged over their components). For each
test composition we computed the Pearson correlation between the
prediction-error time course and each signature. To test whether the full
error vector, rather than its magnitude alone, is linearly decodable from
the population, we fitted ordinary least-squares regressions from
$\mathbf{h}_{t+1}$ (and from $\mathbf{h}_{t+1}-\mathbf{h}_t$) to
$\mathbf{e}_t$ on the training compositions and evaluated the mean squared
decoding error on each test composition, against a baseline decoder that
constantly predicts the training-set mean error vector.

\subsection*{Implementation}

All models were implemented in PyTorch~\cite{paszke2019pytorch}; MIDI
parsing used mido, and statistical analyses used SciPy and
scikit-learn~\cite{pedregosa2011sklearn}. The full pipeline---corpus
retrieval, model training, benchmarking and analysis---is available at
\url{https://github.com/USERNAME/BayesPlusBach}. % TODO: final repo URL
). Requires amsmath (loaded by the template).
% Citation keys are placeholders to be resolved in the project .bib file.
% =========================================================================

\section*{Materials and methods}

Our experimental logic proceeds in four steps. First, we trained recurrent
neural networks (RNNs) exclusively to \emph{decode} (denoise) noisy
observations of a temporally structured stimulus stream---Bach
compositions---without any explicit predictive objective. Second, we asked
whether \emph{prediction emerges} as a by-product of this training: freezing
the network, we probed its hidden state with a linear readout trained to
report the upcoming stimulus. Third, we \emph{benchmarked} the probed
prediction accuracy against a ladder of reference models whose expected
errors are ordered by construction, allowing us to attribute performance to
specific computational capacities. Fourth, we tested whether \emph{prediction
errors are encoded} in the population activity of the trained networks, as
posited by predictive-coding accounts of cortical
processing~\cite{rao1999predictive,friston2005theory}.

\subsection*{Stimulus corpus}

We assembled a corpus of works by J.~S.~Bach from publicly available MIDI
archives.\footnote{jsbach.net, metronimo.com, piano-midi.de, sound.jp/bach-ken,
dardel.info, classicalarchives.com, and suzumidi.com.} Each MIDI file was
parsed into an event-based sequence of chords: a new element was created
whenever the set of simultaneously sounding pitches changed, so that the
sequence index $t$ advances with the musical surface rather than with clock
time. This discretisation removes tempo and rhythm while preserving the
harmonic and voice-leading structure of the music. Pitches were collapsed
into the 12 pitch classes, yielding for each time step a multi-hot binary
vector $\mathbf{y}_t \in \{0,1\}^{12}$ whose $k$-th component indicates
whether pitch class $k$ is sounding (mean polyphony $\approx 2.3$
simultaneous pitch classes). Files belonging to the same composition
(different arrangements or movements sharing a BWV catalogue number) were
concatenated, separated by brief silences.

The corpus comprised 700 compositions, split at the composition level into
training (70\%, 490 works), validation (10\%, 70 works) and test sets (20\%,
140 works). Splitting by BWV number guarantees that no arrangement or
movement of a test composition was ever seen during training.

\subsection*{Noise model and decoding task}

Networks never observed the stimulus directly. At each time step the input
was a noise-corrupted observation
\begin{equation}
  \mathbf{x}_t \;=\; \mathbf{y}_t + \sigma\,\boldsymbol{\varepsilon}_t,
  \qquad \boldsymbol{\varepsilon}_t \sim \mathcal{N}(\mathbf{0}, I_{12}),
  \label{eq:noise}
\end{equation}
with the noise amplitude $\sigma$ treated as an experimental factor spanning
nine levels: a near-noiseless control ($\sigma = 0.001$) and $\sigma \in
\{0.2, 0.4, 0.6, 0.8, 1.0, 1.2, 1.6, 2.0\}$, sampled densely in the range
where the single-observation discriminability $d' = 1/\sigma$ crosses unity
and sparsely near the two asymptotic regimes. At the upper end of this range
the single-step signal-to-noise ratio is far below unity, so that accurate
decoding is impossible without integrating information over time---the
regime in which an internal model of the stimulus dynamics becomes
computationally necessary, in analogy to optimal filtering of a hidden
dynamical process~\cite{kalman1960}. During training,
an independent noise realisation was drawn at every presentation of a
stimulus segment, making the effective training set unbounded and preventing
memorisation of any particular noise pattern. Validation and test sets used
fixed noise realisations, making model selection and evaluation
deterministic.

\subsection*{Network architecture}

The core model was a single-layer gated recurrent unit network
(GRU~\cite{cho2014gru}) with $n$ hidden units, $n \in \{8, 16, 32, 64, 128,
256\}$; the smallest size exceeds the stimulus dimensionality, so no network
was representationally bottlenecked below its input. Two parallel readout heads, each a linear map followed by an
element-wise sigmoid, projected the hidden state $\mathbf{h}_t \in
\mathbb{R}^{n}$ to the stimulus space: an \emph{observation head}
$\hat{\mathbf{y}}_t = \varsigma(W_{o}\mathbf{h}_t + \mathbf{b}_o)$ reporting
the current chord, and a \emph{prediction head} $\mathbf{p}_t =
\varsigma(W_{p}\mathbf{h}_t + \mathbf{b}_p)$ reporting the upcoming chord.
Outputs were treated as independent Bernoulli probabilities over the 12
pitch classes and scored with the binary cross-entropy (BCE),
\begin{equation}
  \mathcal{L}_{\mathrm{obs}}
    = \big\langle \mathrm{BCE}(\hat{\mathbf{y}}_t, \mathbf{y}_t) \big\rangle_t,
  \qquad
  \mathcal{L}_{\mathrm{pred}}
    = \big\langle \mathrm{BCE}(\mathbf{p}_t, \mathbf{y}_{t+1}) \big\rangle_t,
  \label{eq:objectives}
\end{equation}
where $\langle\cdot\rangle_t$ averages over time steps and pitch classes.
The hidden state was initialised to zero at the beginning of each sequence.

\subsection*{Stage 1: training to decode}

In the first training stage only the recurrent weights and the observation
head were optimised, under the decoding objective
$\mathcal{L}_{\mathrm{obs}}$ alone; the prediction head was frozen at its
random initialisation and contributed no gradient. Networks were trained
with Adam~\cite{kingma2015adam} (learning rate $0.02$) on mini-batches of
512 segments of 512 consecutive time steps, drawn by advancing through
randomly selected training compositions. Each network received a fixed
budget of 20{,}000 batches; no early stopping was used, so that every cell
of the $\sigma \times n$ design received an identical amount of training.
Decoding error on the validation set was measured every 100 batches, and the
network state with the lowest validation error was retained as the final
model (best-checkpoint selection). Five networks with independent random
initialisations were trained for each of the $9 \times 6$ combinations of
noise amplitude and network size.

\subsection*{Stage 2: probing for emergent prediction}

To ask whether the decoding-trained networks had learned to anticipate their
input, we froze the recurrent weights and the observation head and trained
\emph{only} the linear prediction head, under
$\mathcal{L}_{\mathrm{pred}}$, for 3{,}000 batches with otherwise identical
hyperparameters. Because the recurrent dynamics were fixed, this readout
cannot create new temporal computations: its accuracy is an estimate of the
predictive information already present in the hidden state. Any prediction
performance exceeding that of the memoryless reference models below must
therefore have emerged spontaneously from the denoising objective of
Stage~1.

\subsection*{Reference models}

Prediction accuracy is only interpretable relative to the intrinsic
predictability of the stimulus. We therefore constructed a ladder of
reference models whose expected losses, at their respective optima,
correspond to conditional entropies of the stimulus process and are
provably ordered. Writing $H(\cdot\,|\,\cdot)$ for the conditional entropy
of a pitch class, independence of the observation noise implies the Markov
chain $\mathbf{x}_t \to \mathbf{y}_t \to \mathbf{y}_{t+1}$, and hence
\begin{equation}
  \underbrace{H(\mathbf{y}_{t+1} \,|\, \mathbf{y}_t)}_{\text{Markov-low}}
  \;\le\;
  \underbrace{H(\mathbf{y}_{t+1} \,|\, \mathbf{x}_t)}_{\text{Markov-high}}
  \;\le\;
  \underbrace{H(\mathbf{y}_{t+1})}_{\text{global distribution}} .
  \label{eq:ladder}
\end{equation}
The corresponding estimators were:
\begin{description}
  \item[Global distribution.] The optimal static predictor under the BCE
  loss: the vector of per-pitch-class marginal frequencies estimated from
  the training set, constantly predicted at every time step. This model has
  no access to temporal structure and upper-bounds the error of any
  informed predictor.
  \item[Markov-high.] A feedforward network (one hidden layer of 256 units)
  mapping the noisy observation $\mathbf{x}_t$ to $\mathbf{y}_{t+1}$: the
  best one-step predictor operating, like the RNN, through the noisy
  channel, but without memory.
  \item[Markov-low.] The same architecture mapping the \emph{clean} chord
  $\mathbf{y}_t$ to $\mathbf{y}_{t+1}$, trained and evaluated with clean
  input: the best one-step predictor with perfect knowledge of the current
  state. Its error is independent of $\sigma$ and lower-bounds all
  first-order predictors.
  \item[Prediction-trained GRU.] A single-layer GRU of 256 units trained
  directly on $\mathcal{L}_{\mathrm{pred}}$, providing an empirical
  estimate of the best prediction attainable by the architecture itself.
  \item[Memoryless decoder.] For the decoding task, a feedforward network
  (three hidden layers of 256 units) mapping $\mathbf{x}_t$ to
  $\mathbf{y}_t$, quantifying how much decoding accuracy is attainable
  without temporal integration.
\end{description}
Because the bound interpretation of Eq~\eqref{eq:ladder} requires reference
models operating near their optima, all reference models were
capacity-saturated (256 hidden units, the maximum of the network-size
sweep) and computed once per noise level, rather than width-matched to each
RNN of the sweep. They were trained with the same optimiser, batch
geometry, validation cadence and checkpoint selection as the core networks,
with a budget of 25{,}000 batches.

\subsection*{Evaluation and statistics}

All models were evaluated on the 140 held-out test compositions, presented
in full and never seen during training or model selection. The test-set
noise realisation was generated once from a fixed seed and shared by all
models and all noise amplitudes (common random numbers: the same
$\boldsymbol{\varepsilon}$ scaled by $\sigma$), so that comparisons between
models are paired both by composition and by noise realisation. Performance
was summarised as the mean BCE per composition; for the core networks,
per-composition errors were first averaged across the five training runs, so
that the composition is the unit of analysis throughout. This yields for
each model a distribution of 140 paired errors. Pairwise comparisons
between models used
two-sided Wilcoxon signed-rank tests across test compositions, with effect
sizes reported as Cohen's $d$ computed with the pooled standard deviation.
For summary maps spanning the full design, $p$-values were
Bonferroni-corrected for the $9 \times 6 = 54$ cells of the sweep.

\subsection*{Encoding of prediction errors}

Finally, we asked whether the trained networks explicitly represent their
own prediction errors, defined at each time step as the deviation of the
incoming observation from the network's prediction,
\begin{equation}
  \mathbf{e}_t \;=\; \mathbf{x}_{t+1} - \mathbf{p}_t .
  \label{eq:pe}
\end{equation}
Note that $\mathbf{e}_t$ is defined with respect to the noisy observation,
the only signal available to the network, so its magnitude scales with
$\sigma$. We related the population prediction error $\|\mathbf{e}_t\|^2$ to
two signatures of the hidden dynamics: the state energy
$\|\mathbf{h}_{t+1}\|^2$ and the state-change energy $\|\mathbf{h}_{t+1} -
\mathbf{h}_t\|^2$ (all quantities averaged over their components). For each
test composition we computed the Pearson correlation between the
prediction-error time course and each signature. To test whether the full
error vector, rather than its magnitude alone, is linearly decodable from
the population, we fitted ordinary least-squares regressions from
$\mathbf{h}_{t+1}$ (and from $\mathbf{h}_{t+1}-\mathbf{h}_t$) to
$\mathbf{e}_t$ on the training compositions and evaluated the mean squared
decoding error on each test composition, against a baseline decoder that
constantly predicts the training-set mean error vector.

\subsection*{Implementation}

All models were implemented in PyTorch~\cite{paszke2019pytorch}; MIDI
parsing used mido, and statistical analyses used SciPy and
scikit-learn~\cite{pedregosa2011sklearn}. The full pipeline---corpus
retrieval, model training, benchmarking and analysis---is available at
\url{https://github.com/USERNAME/BayesPlusBach}. % TODO: final repo URL
). Requires amsmath (loaded by the template).
% Citation keys are placeholders to be resolved in the project .bib file.
% =========================================================================

\section*{Materials and methods}

Our experimental logic proceeds in four steps. First, we trained recurrent
neural networks (RNNs) exclusively to \emph{decode} (denoise) noisy
observations of a temporally structured stimulus stream---Bach
compositions---without any explicit predictive objective. Second, we asked
whether \emph{prediction emerges} as a by-product of this training: freezing
the network, we probed its hidden state with a linear readout trained to
report the upcoming stimulus. Third, we \emph{benchmarked} the probed
prediction accuracy against a ladder of reference models whose expected
errors are ordered by construction, allowing us to attribute performance to
specific computational capacities. Fourth, we tested whether \emph{prediction
errors are encoded} in the population activity of the trained networks, as
posited by predictive-coding accounts of cortical
processing~\cite{rao1999predictive,friston2005theory}.

\subsection*{Stimulus corpus}

We assembled a corpus of works by J.~S.~Bach from publicly available MIDI
archives.\footnote{jsbach.net, metronimo.com, piano-midi.de, sound.jp/bach-ken,
dardel.info, classicalarchives.com, and suzumidi.com.} Each MIDI file was
parsed into an event-based sequence of chords: a new element was created
whenever the set of simultaneously sounding pitches changed, so that the
sequence index $t$ advances with the musical surface rather than with clock
time. This discretisation removes tempo and rhythm while preserving the
harmonic and voice-leading structure of the music. Pitches were collapsed
into the 12 pitch classes, yielding for each time step a multi-hot binary
vector $\mathbf{y}_t \in \{0,1\}^{12}$ whose $k$-th component indicates
whether pitch class $k$ is sounding (mean polyphony $\approx 2.3$
simultaneous pitch classes). Files belonging to the same composition
(different arrangements or movements sharing a BWV catalogue number) were
concatenated, separated by brief silences.

The corpus comprised 700 compositions, split at the composition level into
training (70\%, 490 works), validation (10\%, 70 works) and test sets (20\%,
140 works). Splitting by BWV number guarantees that no arrangement or
movement of a test composition was ever seen during training.

\subsection*{Noise model and decoding task}

Networks never observed the stimulus directly. At each time step the input
was a noise-corrupted observation
\begin{equation}
  \mathbf{x}_t \;=\; \mathbf{y}_t + \sigma\,\boldsymbol{\varepsilon}_t,
  \qquad \boldsymbol{\varepsilon}_t \sim \mathcal{N}(\mathbf{0}, I_{12}),
  \label{eq:noise}
\end{equation}
with the noise amplitude $\sigma$ treated as an experimental factor spanning
nine levels: a near-noiseless control ($\sigma = 0.001$) and $\sigma \in
\{0.2, 0.4, 0.6, 0.8, 1.0, 1.2, 1.6, 2.0\}$, sampled densely in the range
where the single-observation discriminability $d' = 1/\sigma$ crosses unity
and sparsely near the two asymptotic regimes. At the upper end of this range
the single-step signal-to-noise ratio is far below unity, so that accurate
decoding is impossible without integrating information over time---the
regime in which an internal model of the stimulus dynamics becomes
computationally necessary, in analogy to optimal filtering of a hidden
dynamical process~\cite{kalman1960}. During training,
an independent noise realisation was drawn at every presentation of a
stimulus segment, making the effective training set unbounded and preventing
memorisation of any particular noise pattern. Validation and test sets used
fixed noise realisations, making model selection and evaluation
deterministic.

\subsection*{Network architecture}

The core model was a single-layer gated recurrent unit network
(GRU~\cite{cho2014gru}) with $n$ hidden units, $n \in \{8, 16, 32, 64, 128,
256\}$; the smallest size exceeds the stimulus dimensionality, so no network
was representationally bottlenecked below its input. Two parallel readout heads, each a linear map followed by an
element-wise sigmoid, projected the hidden state $\mathbf{h}_t \in
\mathbb{R}^{n}$ to the stimulus space: an \emph{observation head}
$\hat{\mathbf{y}}_t = \varsigma(W_{o}\mathbf{h}_t + \mathbf{b}_o)$ reporting
the current chord, and a \emph{prediction head} $\mathbf{p}_t =
\varsigma(W_{p}\mathbf{h}_t + \mathbf{b}_p)$ reporting the upcoming chord.
Outputs were treated as independent Bernoulli probabilities over the 12
pitch classes and scored with the binary cross-entropy (BCE),
\begin{equation}
  \mathcal{L}_{\mathrm{obs}}
    = \big\langle \mathrm{BCE}(\hat{\mathbf{y}}_t, \mathbf{y}_t) \big\rangle_t,
  \qquad
  \mathcal{L}_{\mathrm{pred}}
    = \big\langle \mathrm{BCE}(\mathbf{p}_t, \mathbf{y}_{t+1}) \big\rangle_t,
  \label{eq:objectives}
\end{equation}
where $\langle\cdot\rangle_t$ averages over time steps and pitch classes.
The hidden state was initialised to zero at the beginning of each sequence.

\subsection*{Stage 1: training to decode}

In the first training stage only the recurrent weights and the observation
head were optimised, under the decoding objective
$\mathcal{L}_{\mathrm{obs}}$ alone; the prediction head was frozen at its
random initialisation and contributed no gradient. Networks were trained
with Adam~\cite{kingma2015adam} (learning rate $0.02$) on mini-batches of
512 segments of 512 consecutive time steps, drawn by advancing through
randomly selected training compositions. Each network received a fixed
budget of 20{,}000 batches; no early stopping was used, so that every cell
of the $\sigma \times n$ design received an identical amount of training.
Decoding error on the validation set was measured every 100 batches, and the
network state with the lowest validation error was retained as the final
model (best-checkpoint selection). Five networks with independent random
initialisations were trained for each of the $9 \times 6$ combinations of
noise amplitude and network size.

\subsection*{Stage 2: probing for emergent prediction}

To ask whether the decoding-trained networks had learned to anticipate their
input, we froze the recurrent weights and the observation head and trained
\emph{only} the linear prediction head, under
$\mathcal{L}_{\mathrm{pred}}$, for 3{,}000 batches with otherwise identical
hyperparameters. Because the recurrent dynamics were fixed, this readout
cannot create new temporal computations: its accuracy is an estimate of the
predictive information already present in the hidden state. Any prediction
performance exceeding that of the memoryless reference models below must
therefore have emerged spontaneously from the denoising objective of
Stage~1.

\subsection*{Reference models}

Prediction accuracy is only interpretable relative to the intrinsic
predictability of the stimulus. We therefore constructed a ladder of
reference models whose expected losses, at their respective optima,
correspond to conditional entropies of the stimulus process and are
provably ordered. Writing $H(\cdot\,|\,\cdot)$ for the conditional entropy
of a pitch class, independence of the observation noise implies the Markov
chain $\mathbf{x}_t \to \mathbf{y}_t \to \mathbf{y}_{t+1}$, and hence
\begin{equation}
  \underbrace{H(\mathbf{y}_{t+1} \,|\, \mathbf{y}_t)}_{\text{Markov-low}}
  \;\le\;
  \underbrace{H(\mathbf{y}_{t+1} \,|\, \mathbf{x}_t)}_{\text{Markov-high}}
  \;\le\;
  \underbrace{H(\mathbf{y}_{t+1})}_{\text{global distribution}} .
  \label{eq:ladder}
\end{equation}
The corresponding estimators were:
\begin{description}
  \item[Global distribution.] The optimal static predictor under the BCE
  loss: the vector of per-pitch-class marginal frequencies estimated from
  the training set, constantly predicted at every time step. This model has
  no access to temporal structure and upper-bounds the error of any
  informed predictor.
  \item[Markov-high.] A feedforward network (one hidden layer of 256 units)
  mapping the noisy observation $\mathbf{x}_t$ to $\mathbf{y}_{t+1}$: the
  best one-step predictor operating, like the RNN, through the noisy
  channel, but without memory.
  \item[Markov-low.] The same architecture mapping the \emph{clean} chord
  $\mathbf{y}_t$ to $\mathbf{y}_{t+1}$, trained and evaluated with clean
  input: the best one-step predictor with perfect knowledge of the current
  state. Its error is independent of $\sigma$ and lower-bounds all
  first-order predictors.
  \item[Prediction-trained GRU.] A single-layer GRU of 256 units trained
  directly on $\mathcal{L}_{\mathrm{pred}}$, providing an empirical
  estimate of the best prediction attainable by the architecture itself.
  \item[Memoryless decoder.] For the decoding task, a feedforward network
  (three hidden layers of 256 units) mapping $\mathbf{x}_t$ to
  $\mathbf{y}_t$, quantifying how much decoding accuracy is attainable
  without temporal integration.
\end{description}
Because the bound interpretation of Eq~\eqref{eq:ladder} requires reference
models operating near their optima, all reference models were
capacity-saturated (256 hidden units, the maximum of the network-size
sweep) and computed once per noise level, rather than width-matched to each
RNN of the sweep. They were trained with the same optimiser, batch
geometry, validation cadence and checkpoint selection as the core networks,
with a budget of 25{,}000 batches.

\subsection*{Evaluation and statistics}

All models were evaluated on the 140 held-out test compositions, presented
in full and never seen during training or model selection. The test-set
noise realisation was generated once from a fixed seed and shared by all
models and all noise amplitudes (common random numbers: the same
$\boldsymbol{\varepsilon}$ scaled by $\sigma$), so that comparisons between
models are paired both by composition and by noise realisation. Performance
was summarised as the mean BCE per composition; for the core networks,
per-composition errors were first averaged across the five training runs, so
that the composition is the unit of analysis throughout. This yields for
each model a distribution of 140 paired errors. Pairwise comparisons
between models used
two-sided Wilcoxon signed-rank tests across test compositions, with effect
sizes reported as Cohen's $d$ computed with the pooled standard deviation.
For summary maps spanning the full design, $p$-values were
Bonferroni-corrected for the $9 \times 6 = 54$ cells of the sweep.

\subsection*{Encoding of prediction errors}

Finally, we asked whether the trained networks explicitly represent their
own prediction errors, defined at each time step as the deviation of the
incoming observation from the network's prediction,
\begin{equation}
  \mathbf{e}_t \;=\; \mathbf{x}_{t+1} - \mathbf{p}_t .
  \label{eq:pe}
\end{equation}
Note that $\mathbf{e}_t$ is defined with respect to the noisy observation,
the only signal available to the network, so its magnitude scales with
$\sigma$. We related the population prediction error $\|\mathbf{e}_t\|^2$ to
two signatures of the hidden dynamics: the state energy
$\|\mathbf{h}_{t+1}\|^2$ and the state-change energy $\|\mathbf{h}_{t+1} -
\mathbf{h}_t\|^2$ (all quantities averaged over their components). For each
test composition we computed the Pearson correlation between the
prediction-error time course and each signature. To test whether the full
error vector, rather than its magnitude alone, is linearly decodable from
the population, we fitted ordinary least-squares regressions from
$\mathbf{h}_{t+1}$ (and from $\mathbf{h}_{t+1}-\mathbf{h}_t$) to
$\mathbf{e}_t$ on the training compositions and evaluated the mean squared
decoding error on each test composition, against a baseline decoder that
constantly predicts the training-set mean error vector.

\subsection*{Implementation}

All models were implemented in PyTorch~\cite{paszke2019pytorch}; MIDI
parsing used mido, and statistical analyses used SciPy and
scikit-learn~\cite{pedregosa2011sklearn}. The full pipeline---corpus
retrieval, model training, benchmarking and analysis---is available at
\url{https://github.com/USERNAME/BayesPlusBach}. % TODO: final repo URL
